% =====================================================================
%  5.3 Sinussetningen — teori og oppgaver
% =====================================================================

\del{1}{Sinussetningen og beviset}

\begin{ordliste}
sinussetningen & Sammenhengen $\frac{\sin A}{a} = \frac{\sin B}{b} = \frac{\sin C}{c}$ i en trekant. & \mittsprak\\
motstående side & Siden som ligger rett overfor en vinkel. Side $a$ er motstående til vinkel $A$. & \mittsprak\\
vinkelsum & Summen av de tre vinklene i en trekant. Den er alltid $180^\circ$. & \mittsprak\\
\end{ordliste}

\begin{mustkunne}
\begin{minipage}[t]{0.58\linewidth}
\textbf{Sinussetningen.} I trekant $ABC$ står $a$ overfor $A$, $b$ overfor $B$ og $c$ overfor $C$. Da er
\[ \frac{\sin A}{a} = \frac{\sin B}{b} = \frac{\sin C}{c} \qquad\text{eller}\qquad \frac{a}{\sin A} = \frac{b}{\sin B} = \frac{c}{\sin C} \]
Hver vinkel hører sammen med siden overfor.\par
\textbf{Slik bruker du den} når to vinkler og én side er kjent: (1) Finn tredje vinkel med $180^\circ$. (2) Bruk et kjent par (vinkel og motstående side) og den ukjente siden.
\end{minipage}\hfill
\begin{minipage}[t]{0.38\linewidth}
\centering
\begin{tikzpicture}[baseline=(current bounding box.north),scale=0.85,font=\small]
  \coordinate (A) at (0,0); \coordinate (B) at (4,0); \coordinate (C) at (1.4,2.3);
  \draw[thick] (A) -- (B) node[midway,below] {$c$};
  \draw[thick] (B) -- (C) node[midway,right] {$a$};
  \draw[thick] (C) -- (A) node[midway,left] {$b$};
  \pic[draw, angle radius=7mm] {angle = B--A--C};
  \pic[draw, angle radius=6mm] {angle = C--B--A};
  \pic[draw, angle radius=5mm] {angle = A--C--B};
  \node[below left] at (A) {$A$}; \node[below right] at (B) {$B$}; \node[above] at (C) {$C$};
\end{tikzpicture}
\end{minipage}\par\medskip
\textbf{Beviset.} Arealsetningen (5.2) gir arealet tre ganger, med to sider og vinkelen mellom dem:
\[ T = \tfrac12\, bc \sin A = \tfrac12\, ac \sin B = \tfrac12\, ab \sin C \]
Del alle tre uttrykkene på $\tfrac12\, abc$. Nesten alt forkortes bort:
\[ \frac{\sin A}{a} = \frac{\sin B}{b} = \frac{\sin C}{c} \]
Arealsetningen gjelder også for stumpe vinkler. Derfor gjelder sinussetningen i alle trekanter.
\end{mustkunne}

\begin{eksempel}[1 — eksakt]
I trekant $ABC$ er $A = 30^\circ$, $B = 45^\circ$ og $a = 4$. Finn $b$ eksakt.

\begin{stegtabell}
$\dfrac{b}{\sin 45^\circ} = \dfrac{4}{\sin 30^\circ}$ & Sinussetningen. $b$ står overfor $B$ og $4$ overfor $A$.\\
$b = \dfrac{4 \cdot \sin 45^\circ}{\sin 30^\circ}$ & Gang med $\sin 45^\circ$ på begge sider.\\
$b = 4 \cdot \dfrac{\sqrt2}{2} \cdot 2$ & Sett inn $\sin 45^\circ = \frac{\sqrt2}{2}$. Å dele på $\sin 30^\circ = \frac12$ er det samme som å gange med $2$.\\
$b = 4\sqrt2$ & Forenkle.\\
Kontroll: & $4\sqrt2 \approx 5,7 > 4$. Vinkel $B$ er større enn $A$, så $b$ skal være større enn $a$. Det stemmer.\\
\end{stegtabell}
\end{eksempel}

\begin{eksempel}[2 — med kalkulator]
I trekant $ABC$ er $A = 40^\circ$, $B = 75^\circ$ og $a = 8,0$. Finn $C$, $b$ og $c$. Svar med én desimal.

\begin{stegtabell}
$C = 180^\circ - 40^\circ - 75^\circ = 65^\circ$ & Vinkelsummen.\\
$b = \dfrac{8,0 \cdot \sin 75^\circ}{\sin 40^\circ}$ & Sinussetningen med paret $A$ og $a$.\\
$b \approx 12,0$ & Kalkulatoren står i DEG.\\
$c = \dfrac{8,0 \cdot \sin 65^\circ}{\sin 40^\circ} \approx 11,3$ & Samme par, nå med $C$.\\
Kontroll: & Den største vinkelen ($75^\circ$) står overfor den lengste siden ($b$). Det stemmer.\\
\end{stegtabell}
\end{eksempel}

\begin{vanligfeil}
\begin{itemize}[leftmargin=1.2em,itemsep=1pt,topsep=0pt]
  \item $\frac{\sin A}{b} = \frac{\sin B}{a}$ er feil. Vinkelen må stå sammen med siden overfor: $\frac{\sin A}{a} = \frac{\sin B}{b}$.
  \item Å sette $C$ inn uten å regne den ut er feil. Finn først $C = 180^\circ - A - B$.
  \item Sinussetningen trenger et kjent par av vinkel og motstående side. Med bare tre sider, eller to sider og vinkelen mellom dem, har vi ikke det. Da trengs cosinussetningen (5.4).
\end{itemize}
\end{vanligfeil}

\nivaahode{1}

\oppg{1.1}{1}{\uh}
Finn den tredje vinkelen i trekant $ABC$.
\begin{tasks}(2)
  \task $A = 50^\circ$ og $B = 70^\circ$
  \task $A = 35^\circ$ og $B = 95^\circ$
  \task $B = 90^\circ$ og $C = 37^\circ$
  \task $A = B = 65^\circ$
\end{tasks}

\oppg{1.2}{1}{\uh}
I trekant $ABC$ er $A = 60^\circ$, $B = 30^\circ$ og $a = 6\sqrt3$. Fullfør.
\begin{tasks}(1)
  \task Fyll inn: $\dfrac{b}{\sin \skrivlinje[15mm]} = \dfrac{6\sqrt3}{\sin \skrivlinje[15mm]}$
  \task Bruk $\sin 30^\circ = \frac12$ og $\sin 60^\circ = \frac{\sqrt3}{2}$ og finn $b$.
\end{tasks}

\oppg{1.3}{1}{\mh}
I trekant $ABC$ er $A = 50^\circ$, $B = 60^\circ$ og $a = 7,0$. Finn $C$, $b$ og $c$. Svar med én desimal.

\nivaahode{2}

\oppg{1.4}{2}{\mh}
Finn den ukjente siden. Svar med én desimal.
\begin{tasks}(1)
  \task $A = 35^\circ$, $B = 80^\circ$ og $a = 12,0$. Finn $b$.
  \task $B = 48^\circ$, $C = 71^\circ$ og $c = 9,5$. Finn $a$.
  \task $A = 112^\circ$, $B = 28^\circ$ og $b = 5,0$. Finn $a$.
\end{tasks}

\oppg{1.5}{2}{\uh}
I trekant $ABC$ er $A = 120^\circ$, $B = 45^\circ$ og $b = 6\sqrt2$. Finn $a$ eksakt og $C$.

\oppg{1.6}{2}{\mh}
To observatører, $P$ og $Q$, står $500$ m fra hverandre på flat mark. En drone $D$ svever over bakken mellom dem, sett fra siden.
Vinkelen opp til dronen er $62^\circ$ fra $P$ og $71^\circ$ fra $Q$. Tegn en skisse.
\begin{tasks}(1)
  \task Finn vinkelen ved $D$.
  \task Finn avstanden fra $P$ til dronen og fra $Q$ til dronen. Svar med én desimal.
  \task Finn høyden til dronen over bakken. Svar i hele meter.
\end{tasks}

\oppg{1.7}{2}{\mh}
I trekant $ABC$ er $A = 40^\circ$, $B = 65^\circ$ og $a = 9,0$. Eleven skal finne $c$. Finn \textbf{den første feilen}. Løs riktig.
\[ C = 75^\circ \quad\longrightarrow\quad \frac{\sin 75^\circ}{c} = \frac{\sin 40^\circ}{9,0} \quad\longrightarrow\quad c = \frac{9,0 \cdot \sin 40^\circ}{\sin 75^\circ} \quad\longrightarrow\quad c \approx 6,0 \]

\nivaahode{3}

\oppg{1.8}{3}{\uh}
Er påstanden alltid, noen ganger eller aldri sann? Begrunn svaret.
\begin{tasks}(1)
  \task Er $A = B$ i en trekant, er $a = b$.
  \task Sinussetningen gjelder bare i trekanter der alle vinklene er spisse.
  \task Kjenner vi to vinkler og siden mellom dem, kan vi finne de to andre sidene med sinussetningen.
\end{tasks}

\oppg{1.9}{3}{\uh}
I en trekant er $A = 90^\circ$. Sett inn i sinussetningen og vis at $\sin B = \dfrac{b}{a}$.
Forklar hvorfor dette stemmer med definisjonen av sinus i 5.1.

\oppg{1.10}{3}{\uh}
Lag selv en trekant der $b = \sqrt3 \cdot a$. Oppgi $A$, $B$ og $C$, og vis at sinussetningen gir riktig forhold mellom sidene.
Flere svar er riktige.

\begin{sjekk}
\begin{tasks}(2)
  \task Finn $C$ når $A = 75^\circ$ og $B = 45^\circ$.
  \task Finn $b$ når $A = 30^\circ$, $B = 90^\circ$ og $a = 3$.
  \task! Hvorfor må vi finne den tredje vinkelen først når vi kjenner to vinkler og én side?
\end{tasks}
\end{sjekk}

% ---------------------------------------------------------------------

\del{2}{Det tvetydige tilfellet}

\begin{ordliste}
SSV & To sider og en vinkel som \emph{ikke} ligger mellom dem. Vi kjenner for eksempel $A$, $a$ og $b$. & \mittsprak\\
tvetydig & Kan tolkes på to måter. Her: opplysningene passer med to ulike trekanter. & \mittsprak\\
svingarm & Siden $a$ tenkt som en arm som festes i $C$ og svinger ned mot grunnlinja. & \mittsprak\\
\end{ordliste}

\begin{mustkunne}
\begin{minipage}[t]{0.52\linewidth}
\textbf{SSV.} Vi kjenner $A$, $a$ og $b$. Sinussetningen gir
\[ \sin B = \frac{b \sin A}{a} = k \]
Da finnes to vinkler mellom $0^\circ$ og $180^\circ$ med $\sin B = k$ (5.1): $B$ og $180^\circ - B$. Begge er gyldige bare hvis vinkelsummen er under $180^\circ$:
\[ A + B < 180^\circ \]
\textbf{Svingarmen.} $A$ og $b$ ligger fast. Avstanden fra $C$ til grunnlinja er $h = b \sin A$. Armen $a$ svinger fra $C$.
\end{minipage}\hfill
\begin{minipage}[t]{0.44\linewidth}
\centering
\begin{tikzpicture}[baseline=(current bounding box.north),scale=1.2,font=\small]
  \coordinate (A) at (0,0); \coordinate (C) at (2.77,1.6);
  \coordinate (B1) at (1.57,0); \coordinate (B2) at (3.97,0); \coordinate (H) at (2.77,0);
  \draw (A) -- (4.8,0);
  \draw[thick] (A) -- (C) node[midway,above left] {$b$};
  \draw[thick] (C) -- (B1) node[pos=0.35,left] {$a$};
  \draw[thick] (C) -- (B2) node[pos=0.35,right] {$a$};
  \draw[dotted, thick] (C) ++(220:2) arc (220:320:2);
  \draw[dashed] (C) -- (H);
  \draw (2.57,0)--(2.57,0.2)--(2.77,0.2);
  \node[right] at (2.8,0.45) {$h$};
  \pic[draw, angle radius=7mm] {angle = H--A--C};
  \node[left] at (A) {$A$}; \node[above] at (C) {$C$};
  \node[below left] at (B1) {$B_1$}; \node[below right] at (B2) {$B_2$};
\end{tikzpicture}
\end{minipage}\par\medskip
\renewcommand{\arraystretch}{1.4}
\begin{tabular}{@{}l l l@{}}
$a < h$ & ($k > 1$) & ingen trekant, armen når ikke ned.\\
$a = h$ & ($k = 1$) & én trekant, $B = 90^\circ$.\\
$h < a < b$ & ($k < 1$) & to trekanter: $B$ og $180^\circ - B$.\\
$a \ge b$ & & én trekant. Da gir $180^\circ - B$ sum $\ge 180^\circ$.\\
\end{tabular}
\end{mustkunne}

\begin{eksempel}[— to trekanter]
I trekant $ABC$ er $A = 30^\circ$, $a = 5$ og $b = 8$. Finn alle mulige trekanter.

\begin{stegtabell}
$\sin B = \dfrac{8 \cdot \sin 30^\circ}{5} = 0,8$ & Sinussetningen. $h = 8 \sin 30^\circ = 4 < 5 < 8$, så vi venter to trekanter.\\
$B_1 \approx 53,1^\circ$ og $B_2 \approx 126,9^\circ$ & $\sin^{-1}(0,8)$ og $180^\circ - 53,1^\circ$.\\
$30^\circ + 126,9^\circ = 156,9^\circ < 180^\circ$ & Vinkelsummen tillater begge.\\
$C_1 \approx 96,9^\circ$ og $C_2 \approx 23,1^\circ$ & $C = 180^\circ - 30^\circ - B$.\\
$c_1 = \dfrac{5 \sin 96,9^\circ}{\sin 30^\circ} \approx 9,9$ & Sinussetningen med paret $A$ og $a$.\\
$c_2 = \dfrac{5 \sin 23,1^\circ}{\sin 30^\circ} \approx 3,9$ & Samme regning for den andre trekanten.\\
\end{stegtabell}
Kontroll: $c_1 = 4\sqrt3 + 3 \approx 9,9$ og $c_2 = 4\sqrt3 - 3 \approx 3,9$ langs grunnlinja i figuren.
\end{eksempel}

\begin{vanligfeil}
\begin{itemize}[leftmargin=1.2em,itemsep=1pt,topsep=0pt]
  \item Å regne ut bare $B = \sin^{-1}(k)$ er feil når $a < b$. Kalkulatoren gir bare den spisse vinkelen. Den stumpe er $180^\circ - B$.
  \item Å godta begge vinklene uten å sjekke er feil. Er $a \ge b$, blir $A + (180^\circ - B) \ge 180^\circ$, og da finnes bare én trekant.
  \item $\sin B > 1$ betyr ikke at du har regnet feil. Det betyr at det ikke finnes noen trekant.
\end{itemize}
\end{vanligfeil}

\nivaahode{1}

\oppg{2.1}{1}{\uh}
Finn alle vinkler $v$ mellom $0^\circ$ og $180^\circ$ med
\begin{tasks}(4)
  \task $\sin v = \frac12$
  \task $\sin v = 1$
  \task $\sin v = 2$
  \task $\sin v = \frac{\sqrt3}{2}$
\end{tasks}

\oppg{2.2}{1}{\uh}
I trekant $ABC$ er $A = 30^\circ$, $a = 6$ og $b = 6\sqrt3$. Fullfør.
\begin{tasks}(1)
  \task Sinussetningen gir $\sin B = \dfrac{6\sqrt3 \cdot \sin 30^\circ}{6}$. Regn ut $\sin B$ eksakt.
  \task Finn de to mulige vinklene $B$.
  \task Sjekk at begge gir en trekant. Finn $C$ og $c$ i hver trekant.
\end{tasks}

\oppg{2.3}{1}{\mh}
I trekant $ABC$ er $A = 40^\circ$, $a = 6,0$ og $b = 7,0$.
\begin{tasks}(1)
  \task Finn $\sin B$ med to desimaler.
  \task Finn begge mulige vinkler $B$. Svar med én desimal.
  \task Gir begge vinklene en trekant? Vis vinkelsummen.
\end{tasks}

\nivaahode{2}

\oppg{2.4}{2}{\mh}
I trekant $ABC$ er $A = 35^\circ$, $a = 7,0$ og $b = 9,0$. Finn alle trekanter som passer.
Finn $B$, $C$ og $c$ i hver. Svar med én desimal.

\oppg{2.5}{2}{\mh}
Hvor mange trekanter $ABC$ finnes? Finn $B$ der det finnes trekanter.
\begin{tasks}(1)
  \task $A = 62^\circ$, $a = 9,0$ og $b = 6,0$. Finn også $C$ og $c$.
  \task $A = 40^\circ$, $a = 3,0$ og $b = 7,0$
  \task $A = 110^\circ$, $a = 10$ og $b = 7,0$. Finn også $C$ og $c$.
\end{tasks}

\oppg{2.6}{2}{\mh}
Fra en havn $A$ går en rett kystlinje. Et fyrtårn $C$ står $8,0$ km fra $A$, og linja $AC$ danner $25^\circ$ med kysten.
En båt $B$ ligger ved kysten og er $6,0$ km fra fyrtårnet.
\begin{tasks}(1)
  \task Tegn en skisse som viser at det er to mulige plasseringer av båten.
  \task Finn avstanden fra havna til båten i begge tilfellene. Svar med én desimal.
\end{tasks}

\oppg{2.7}{2}{\mh}
Eleven skal finne alle trekanter $ABC$ med $A = 48^\circ$, $a = 7,0$ og $b = 9,0$. Finn \textbf{den første feilen}. Løs riktig.
\[ \frac{\sin B}{9,0} = \frac{\sin 48^\circ}{7,0} \;\longrightarrow\; \sin B \approx 0,955 \;\longrightarrow\; B \approx 72,8^\circ \text{ (eneste løsning)} \;\longrightarrow\; C \approx 59,2^\circ \;\longrightarrow\; c \approx 8,1 \]

\nivaahode{3}

\oppg{2.8}{3}{\mh}
I en trekant er $a = 6,0$ og $b = 9,0$. Vinkelen $A$ kan variere.
For hvilke verdier av $A$ finnes det to trekanter, én trekant og ingen trekant?
Begrunn svaret med $\sin B$ og vinkelsummen.

\oppg{2.9}{3}{\uh}
Er påstanden alltid, noen ganger eller aldri sann? Begrunn svaret.
\begin{tasks}(1)
  \task Er $A$ stump, finnes det høyst én trekant med gitt $A$, $a$ og $b$.
  \task Er $a < b$, finnes det to trekanter med gitt $A$, $a$ og $b$.
\end{tasks}

\oppg{2.10}{3}{\uh}
Velg selv en vinkel $A$ og to sider $a$ og $b$.
Lag ett sett som gir nøyaktig to trekanter og ett som gir nøyaktig én trekant.
Vis med regning at det stemmer. Flere svar er riktige.

\begin{sjekk}
\begin{tasks}(2)
  \task Finn begge vinklene $v$ mellom $0^\circ$ og $180^\circ$ med $\sin v = 0,6$.
  \task Hvor mange trekanter finnes når $A = 30^\circ$, $b = 6$ og $a = 2$?
  \task! Hvorfor må vi sjekke at $A + (180^\circ - B) < 180^\circ$ før vi godtar den stumpe vinkelen?
\end{tasks}
\end{sjekk}
